% !TeX root = ../main.tex
\section{Axiomatic semantics}
\label{sec:axiomatic}

This section goes beyond the slides, which only name axiomatic semantics as
the third flavor; the pre- and postconditions of dynamic analysis, in the
semantic analysis lecture, are its starting point. Axiomatic semantics gives
meaning to programs through logical assertions about what is true before and
after they run.

\subsection{Hoare triples}

\begin{definition}
A \textbf{Hoare triple} $\{P\}\ C\ \{Q\}$ says that if the precondition $P$ holds
before $C$ runs and $C$ terminates, then the postcondition $Q$ holds after it.
\end{definition}

For example $\{x = 5\}\ x := x + 1\ \{x = 6\}$, and the triple of the semantic
analysis lecture, $\{x \le 10\}\ \mathtt{while}\ x < 10\ \mathtt{do}\ x := x + 1\ \{x = 10\}$.

\subsection{Rules}

\begin{mathpar}
\inferrule*[right=Assign]{ }{\{P[x := e]\}\ x := e\ \{P\}}
\and
\inferrule*[right=Seq]{\{P\}\ C_1\ \{R\} \\ \{R\}\ C_2\ \{Q\}}{\{P\}\ C_1 ; C_2\ \{Q\}}
\and
\inferrule*[right=If]{\{P \wedge b\}\ C_1\ \{Q\} \\ \{P \wedge \neg b\}\ C_2\ \{Q\}}{\{P\}\ \mathtt{if}\ b\ \mathtt{then}\ C_1\ \mathtt{else}\ C_2\ \{Q\}}
\and
\inferrule*[right=While]{\{I \wedge b\}\ C\ \{I\}}{\{I\}\ \mathtt{while}\ b\ \mathtt{do}\ C\ \{I \wedge \neg b\}}
\and
\inferrule*[right=Cons]{P \Rightarrow P' \\ \{P'\}\ C\ \{Q'\} \\ Q' \Rightarrow Q}{\{P\}\ C\ \{Q\}}
\end{mathpar}
The assignment rule works backward: to have $P$ after $x := e$, $P$ with $e$ in
place of $x$ must hold before. In \textsc{While}, $I$ is the \textbf{loop
invariant}: it holds before the loop, every iteration preserves it, and when
the loop ends $I \wedge \neg b$ must give the postcondition. \textsc{Cons}
strengthens preconditions and weakens postconditions to connect rules.

\paragraph{Partial and total correctness.} $\{P\}\ C\ \{Q\}$ is partial
correctness: it says nothing about termination. Total correctness, written
$[P]\ C\ [Q]$, also requires $C$ to terminate, usually proved with a
\textbf{variant}, a quantity bounded below that decreases at each iteration.

\paragraph{Weakest preconditions.} Instead of guessing a precondition, the
weakest one can be computed:
\[
\begin{array}{l}
\mathit{wp}(x := e, Q) = Q[x := e] \qquad
\mathit{wp}(C_1 ; C_2, Q) = \mathit{wp}(C_1, \mathit{wp}(C_2, Q)) \\
\mathit{wp}(\mathtt{if}\ b\ \mathtt{then}\ C_1\ \mathtt{else}\ C_2, Q) = (b \Rightarrow \mathit{wp}(C_1, Q)) \wedge (\neg b \Rightarrow \mathit{wp}(C_2, Q))
\end{array}
\]
For loops there is no mechanical formula; an invariant is needed.

\subsection{Example: the factorial}

The factorial of Section~\ref{sec:operational}, with $\mathit{fac}$ starting at 1
and $l$ at $n$, has the invariant
\[
I : \mathit{fac} \cdot l! = n! \wedge l \ge 0.
\]
Before the loop, $1 \cdot n! = n!$. An iteration with $l > 0$ multiplies
$\mathit{fac}$ by $l$ and decrements $l$, and
$(\mathit{fac} \cdot l) \cdot (l - 1)! = \mathit{fac} \cdot l!$, so $I$ still holds. When
the loop ends, $l = 0$ and $I$ gives $\mathit{fac} = n!$. The variant is $l$.
Figure~\ref{fig:invariant} places the assertions on the loop.

\begin{figure}[h]
\centering
\begin{tikzpicture}[box/.style={draw, minimum height=7mm, font=\small, align=center}, cond/.style={draw, diamond, aspect=2.2, font=\small, inner sep=1pt}]
  \node[box] (init) at (0,0) {$\mathit{fac} := 1;\ l := n$};
  \node[cond] (test) at (0,-1.7) {$l > 0$};
  \node[box] (body) at (0,-3.6) {$\mathit{fac} := \mathit{fac} * l;\ l := l - 1$};
  \node[box] (end) at (4.6,-1.7) {$\mathit{fac} = n!$};
  \draw[flecha] (init) -- (test);
  \draw[flecha] (test) -- node[left, font=\footnotesize] {true} (body);
  \draw[flecha] (body.west) -- ++(-1,0) |- (test.west);
  \draw[flecha] (test) -- node[above, font=\footnotesize] {false} (end);
  \node[font=\footnotesize, anchor=west] at (0.2,-0.8) {$I$};
  \node[font=\footnotesize, anchor=west] at (0.2,-2.7) {$I \wedge l > 0$};
  \node[font=\footnotesize, anchor=north] at (4.6,-2.2) {$I \wedge l = 0$};
\end{tikzpicture}
\caption{The invariant $I$ on the factorial loop.}
\label{fig:invariant}
\end{figure}

For a loop that adds $a[0..n-1]$ into $s$, the invariant is
$s = \sum_{k=0}^{i-1} a[k] \wedge 0 \le i \le n$; at the end $i = n$ and $s$ is the full sum.

\subsection{Exercises}

\begin{exercise}
Prove $\{x = n\}\ y := x * x\ \{y = n^2\}$.
\end{exercise}

\paragraph{Solution.} By \textsc{Assign}, $\{x * x = n^2\}\ y := x * x\ \{y = n^2\}$.
Since $x = n \Rightarrow x * x = n^2$, \textsc{Cons} gives the triple.

\begin{exercise}
Compute $\mathit{wp}(x := x + 1, x > 0)$.
\end{exercise}

\paragraph{Solution.} $(x > 0)[x := x + 1] = x + 1 > 0$, that is, $x \ge 0$ on integers.

\begin{exercise}
Is $\{x > 0\}\ \mathtt{while}\ x \ne 0\ \mathtt{do}\ x := x - 1\ \{x = 0\}$ partial or total
correctness? Which variant proves termination?
\end{exercise}

\paragraph{Solution.} With invariant $x \ge 0$ it holds as partial correctness:
at the end $x \ge 0 \wedge x = 0$. It also terminates: $x$ is a variant, bounded
below by 0 and decreasing by one in each iteration.
